\documentclass[10pt]{mypackage}

\usepackage[uprightgreeks,light]{kpfonts}
\usepackage{homework}
%\usepackage{notes}

%\usepackage[ backend=bibtex, style = alphabetic, sorting=ynt ]{biblatex}
%\addbibresource{  }

\usepackage{parskip}

\fancyhf{}
\fancyhead[R]{Avinash Iyer}
\fancyhead[L]{Algebraic Topology: Homework 21}
\fancyfoot[C]{\thepage}

\setcounter{secnumdepth}{0}

\begin{document}
\RaggedRight
\begin{problem}[Problem 1]
  Let $\left( X,A \right)$ be a pair of spaces. 
  \begin{enumerate}[(a)]
    \item Prove that $H_0\left( X,A \right) = 0$ if and only if $A$ intersects each path component of $X$ non-trivially.
    \item Prove that $H_1\left( X,A \right) = 0$ if and only if the map $H_1(A)\rightarrow H_1(X)$ is surjective and each path component of $X$ contains at most one path component of $A$.
  \end{enumerate}
\end{problem}
\begin{solution}\hfill
  \begin{enumerate}[(a)]
    \item Suppose $H_0\left( X,A \right) = 0$. Then, we have that $H_0\left( X \right)\cong H_0\left( A \right)$; since $H_0\left( X \right)$ corresponds to path components in $X$ and $H_0\left( A \right)$ corresponds to path components in $A$, it follows that each path component of $X$ includes a path component of $A$, meaning that $A$ intersects each path component of $X$ non-trivially.

      Now, if $A$ intersects each path component of $X$ non-trivially, then there is an injection from the path components of $X$ into the path components of $A$, meaning that $H_0\left( X \right)\leq H_0\left( A \right)$ as subgroups. Yet, since the boundary map on chains descends to the quotients $C_0\left( X \right)/C_0\left( A \right)$ and $C_1\left( X \right)/C_1\left( A \right)$, we must have $H_0\left( A \right)\leq H_0\left( X \right)$ are subgroups, so that $H_0\left( X \right)\cong H_0\left( A \right)$, and thus $H_0\left( X,A \right) = 0$.
    \item  In the long exact sequence below, we observe that if $H_1\left( X,A \right) = 0$, then the corresponding parts of the sequence are given by
      \begin{align*}
        H_1\left( A \right)\rightarrow H_1\left( X \right)\rightarrow 0\\
        0 \rightarrow H_0\left( A \right)\rightarrow H_0\left( X \right),
      \end{align*}
      so by exactness, we have that $H_1\left( A \right)$ surjects onto $H_1\left( X \right)$, and $H_0\left( A \right)$ injects into $H_0\left( X \right)$. The latter implies that the path components of $A$ are contained in the path components of $X$ as $H_0$ is indexed by path components.
      \begin{center}
        % https://tikzcd.yichuanshen.de/#N4Igdg9gJgpgziAXAbVABwnAlgFyxMJZARgBoAGAXVJADcBDAGwFcYkQAJAfWIAoBBAJQgAvqXSZc+QigBMFanSat23PgA1hYidjwEiAZgU0GLNok49e60kNHiQGXdKIAWY0rOqu5AVodOUvooAKwepioW3L6a9jpBMsjk4crmIAA66QDGUBA4CNqOknqJAGwpXhaZOXkFijBQAObwRKAAZgBOEAC2SGEgOBBI5IWdPcM0g0hknpEZ6fg49FzAmfRwOCJxIGO9iDNTiPKzaZloWCtrG1ujXXvHh0Yn7Gf0HXhM27tIT4fuz1UFnllqt0utNl87kh-odSiJKCIgA
\begin{tikzcd}
\cdots \arrow[r] & H_1(A) \arrow[r, "\iota_{\ast}"] & H_1(X) \arrow[r, "\pi_{\ast}"] & {H_1(X,A)} \arrow[r, "\partial"] & H_0(A) \arrow[r, "\iota_{\ast}"] & H_0(X) \arrow[r] & \cdots
\end{tikzcd}
      \end{center}
      In the reverse direction, we know that if each path component of $X$ contains at most one path component of $A$, then the map $\iota_{\ast}\colon H_0\left( A \right)\rightarrow H_0\left( X \right)$ is injective. Similarly, if the map $H_1\left( A \right)\rightarrow H_1\left( X \right)$. Since $\partial$ preserves exactness, it follows that the entire sequence
      \begin{align*}
        H_1\left( A \right)\rightarrow H_1\left( X \right)\rightarrow H_1\left( X,A \right)\rightarrow H_0\left( A \right)\rightarrow H_0\left( X \right)
      \end{align*}
      is exact, which can only occur if $H_1\left( X,A \right) = 0$.
  \end{enumerate}
\end{solution}
\end{document}
